% SGA TRANSLATION PROJECT
% Volume: SGA 4
% Expose: I, Presheaves
% Source: Orgogozo/Laszlo TeX snapshot, SGA4-master-71766d9/01/01.tex
% Range translated in this working draft section: Section 9 through Lemma 9.13.1, source lines 5805--6359.
% Status: draft translation, compile-checked, not proofed against scan.

\documentclass[11pt]{article}
\usepackage[T1]{fontenc}
\usepackage[utf8]{inputenc}
\usepackage[english]{babel}
\usepackage{lmodern}
\usepackage{microtype}
\usepackage{amsmath,amssymb,amsthm,mathtools,amscd}
\usepackage{tikz-cd}
\usepackage{enumitem}
\usepackage{geometry}
\geometry{margin=1.12in}
\usepackage{hyperref}
\hypersetup{colorlinks=true,linkcolor=blue,citecolor=blue,urlcolor=blue}

\newcommand{\U}{\mathcal{U}}
\newcommand{\V}{\mathcal{V}}
\newcommand{\Set}{\mathbf{Set}}
\newcommand{\Ens}{\mathbf{Set}}
\newcommand{\Ob}{\operatorname{Ob}}
\newcommand{\ob}{\operatorname{ob}}
\newcommand{\Hom}{\operatorname{Hom}}
\newcommand{\Fun}{\operatorname{Fun}}
\newcommand{\Ind}{\operatorname{Ind}}
\newcommand{\Filt}{\operatorname{Filt}}
\newcommand{\Ker}{\operatorname{Ker}}
\newcommand{\Coker}{\operatorname{Coker}}
\newcommand{\Sup}{\operatorname{Sup}}
\newcommand{\card}{\operatorname{card}}
\newcommand{\Fl}{\operatorname{Fl}}
\newcommand{\Hhom}{\operatorname{Hom}}
\newcommand{\acc}{\mathrm{acc}}
\newcommand{\preacc}{\mathrm{preacc}}
\newcommand{\id}{\operatorname{id}}
\newcommand{\cf}{cf.}
\newcommand{\ie}{i.e.}
\newcommand{\resp}{respectively}
\newcommand{\qedtext}{\hfill\(\square\)}
\newcommand{\fprod}[3]{#1\times_{#3}#2}
\newcommand{\emptyE}{\varnothing_E}

\newenvironment{statement}[1]{\par\medskip\noindent\textbf{#1.}\ }{\par\medskip}
\newenvironment{prooftext}{\par\noindent\emph{Proof.}\ }{\par\medskip}
\newenvironment{remarktext}[1]{\par\medskip\noindent\textbf{#1.}\ }{\par\medskip}

\begin{document}

\begin{center}
{\Large \textbf{Expos\'e I. Presheaves}}\\[0.35em]
{\large Section 9, beginning through Lemma 9.13.1}\\[0.35em]
A. Grothendieck and J.-L. Verdier
\end{center}

\bigskip


\setcounter{section}{8}
\section{Accessible functors, cardinal filtrations, and construction of small generating subcategories}

The present section is more technical than the other sections of this expos\'e. In this seminar it will be used only in IV~9 and VI~4, which themselves are not used elsewhere in the Seminar. It is therefore advisable to omit the present section on a first reading.

\subsection*{9.0}
All categories considered in this number are assumed to be $\U$-categories. Except for the small indexing categories $I,J,\ldots$ that we shall use, the developments below will mainly concern ``large'' categories $E,F,\ldots$ that are stable under small filtered inductive limits. Most often, however, it will be enough that a slightly weaker condition hold, namely condition $L$ of Definition~9.1 below. All cardinals considered in this number are assumed to belong to $\U$.

Following a suggestion of P.~Deligne, we shall study, for a functor $f:E\to F$ between large categories, a condition expressing commutation of $f$ with certain types of filtered inductive limits. This condition is remarkably stable and is satisfied by the most important functors that occur naturally. The applications we have in view for this seminar are 9.13.3 and 9.13.4, used in VI~4, and above all 9.25, which gives, in a non-trivial case, the existence of a small generating family in a category of sections of a fibred category; that result will be used in IV~9.16.

\begin{statement}{Definition 9.1}
\begin{enumerate}[label=\alph*)]
\item Let $I$ be a preordered set and let $\pi$ be a cardinal. We say that $I$ is \emph{large with respect to} $\pi$ if $I$ is filtered and if every subset of $I$ of cardinal $\leq \pi$ has an upper bound in $I$.

\item Let $E$ be a category. If $\pi$ is a cardinal, we say that $E$ satisfies condition $L_\pi$ if, for every small ordered set $I$ large with respect to $\pi$, $E$ is stable under inductive limits of type $I$. We say that $E$ satisfies condition $L$ if there exists a cardinal $\pi\in\U$ such that $E$ satisfies $L_\pi$.
\end{enumerate}
\end{statement}

\subsection*{9.1.1}
When, in Definition~9.1(a), one has $\pi\geq 2$, the second stated condition already implies that $I$ is filtered. If $\pi$ is finite, saying that $I$ is large with respect to $\pi$ simply means that $I$ is filtered. In what follows we shall be interested almost exclusively in the case where $\pi$ is infinite. Notice that if $\pi$ and $\pi'$ are cardinals with $\pi'\geq \pi$, then $I$ large with respect to $\pi'$ plainly implies $I$ large with respect to $\pi$.

\subsection*{9.1.2}
As announced in 9.0, the conditions $L_\pi$ and $L$ should be regarded as technical variants of the stronger condition of stability under small filtered inductive limits. It is clear that if $\pi$ and $\pi'$ are cardinals with $\pi'\geq \pi$, then condition $L_\pi$ implies condition $L_{\pi'}$.

\begin{statement}{Definition 9.2}
Let $f:E\to F$ be a functor. If $\Pi$ is a cardinal, we say that $f$ is $\Pi$-\emph{preaccessible} (respectively $\Pi$-\emph{accessible}) if $E$ satisfies $L_\Pi$ (9.1) and if, for every ordered set $I\in\U$ large with respect to $\Pi$, and every inductive system $(X_i)_{i\in I}$ in $E$ of type $I$, the canonical morphism
\[
 \varinjlim f(X_i)\longrightarrow f(\varinjlim X_i)
\]
is a monomorphism (respectively an isomorphism). We say that $f$ is \emph{preaccessible} (respectively \emph{accessible}), relative to the universe $\U$, if there exists a cardinal $\Pi\in\U$ such that $f$ is $\Pi$-preaccessible (respectively $\Pi$-accessible).
\end{statement}

The category of $\Pi$-accessible (respectively accessible) functors from $E$ to $F$ will be denoted by $\Hhom(E,F)_\Pi$ (respectively $\Hhom(E,F)_{\acc}$).

\subsection*{9.2.1}
Plainly, a functor which commutes with small filtered inductive limits, for example a functor admitting a right adjoint, is $\Pi$-accessible for every cardinal $\Pi\geq 2$.

\begin{statement}{Definition 9.3}
Let $E$ be a category, let $X$ be an object of $E$, and let
\[
 h_X^\circ:E\longrightarrow (\U\text{-}\Set)
\]
be the covariant functor it represents. Let $\Pi\in\U$ be a cardinal. We say that $X$ is a $\Pi$-preaccessible (respectively $\Pi$-accessible) object of $E$ if the functor $h_X^\circ$ is $\Pi$-preaccessible (respectively $\Pi$-accessible). We say that $X$ is preaccessible (respectively accessible) if there exists a cardinal $\Pi\in\U$ such that $X$ is a $\Pi$-preaccessible (respectively $\Pi$-accessible) object of $E$.
\end{statement}

\subsection*{9.3.1}
We denote by $E_\Pi$ the strictly full subcategory of $E$ consisting of the $\Pi$-accessible objects of $E$.

When Definition~9.3 is applied to a category of the form $\Hhom(C,F)$, the terminology just introduced is, \emph{a priori}, ambiguous with the analogous terminology introduced in 9.2, when one interprets the objects of $\Hhom(C,F)$ as functors. There does not, however, seem to be any serious risk of confusion.

\begin{statement}{Definition 9.4}
Let $E$ be a category and let $\pi\in\U$ be a cardinal. We say that $E$ is a $\pi$-preaccessible (respectively $\pi$-accessible) category if there exists in $E$ a small full subcategory $C$ which is generating (7.1) and whose objects are $\pi$-preaccessible (respectively $\pi$-accessible) (9.3). We say that $E$ is preaccessible (respectively accessible) if there exists a cardinal $\pi\in\U$ such that $E$ is $\pi$-preaccessible (respectively $\pi$-accessible).
\end{statement}

For important examples, see 9.11.3 below.

\begin{statement}{Proposition 9.5}
Let $f:E\to F$ be a functor between categories such that $F$ is accessible and every object of $E$ is accessible. If $f$ admits a left adjoint, then $f$ is accessible.
\end{statement}

Indeed, by hypothesis, $F$ admits a small conservative family of representable functors $F\to(\U\text{-}\Set)$ which are accessible. Thus one is immediately reduced to showing that the composite of $f$ with each of these functors is accessible. Since $f$ admits a left adjoint, these composites are representable functors $E\to(\U\text{-}\Set)$, hence are accessible by the hypothesis on $E$.

\begin{statement}{Proposition 9.6}
Let $E$ and $F$ be two categories, let $\Pi\in\U$ be a cardinal, and consider the full subcategory $\Hhom(E,F)_\Pi$ of $\Hhom(E,F)$ formed by the $\Pi$-accessible functors (9.2) from $E$ to $F$.
\begin{enumerate}[label=(\roman*)]
\item This subcategory is stable under every type of inductive limit which is representable in $F$.

\item Suppose that $F$ is $\Pi$-accessible (9.4), and let $J$ be a category such that $\card\Fl J\leq \Pi$ and such that projective limits of type $J$ are representable in $F$; then projective limits of type $J$ are representable in $\Hhom(E,F)$. The subcategory $\Hhom(E,F)_\Pi$ is stable under projective limits of type $J$.
\end{enumerate}
\end{statement}

\begin{statement}{Corollary 9.7}
Let $E$ and $F$ be two categories, with $F$ accessible (9.4). Then the full subcategory $\Hhom(E,F)_{\acc}$ of $\Hhom(E,F)$ formed by accessible functors is stable under every inductive or projective limit, relative to a small indexing category $J$, which is representable in $F$ and hence in $\Hhom(E,F)$.
\end{statement}

\begin{prooftext}[Proof of Proposition 9.6]
Assertion (i) follows trivially from the commutation of the functor $\varinjlim$ with arbitrary inductive limits. For (ii), let $(f_j)_{j\in J}$ be a projective system of $\Pi$-accessible functors $E\to F$, and let
\[
 f=\varprojlim_j f_j
\]
be its projective limit, computed argument by argument. We prove that $f$ is $\Pi$-accessible; in other words, for every ordered set $I\in\U$ large with respect to $\Pi$, and every inductive system $(X_i)_{i\in I}$ in $E$, the canonical morphism
\[
 \varinjlim_i ((\varprojlim_j f_j)(X_i))\longrightarrow
 (\varprojlim_j f_j)(\varinjlim_i X_i)
\]
is an isomorphism. The argument-by-argument computation of the functor $\varprojlim_j f_j$ identifies this canonical morphism with
\[
 \varinjlim_i \varprojlim_j f_j(X_i)
 \longrightarrow
 \varprojlim_j \varinjlim_i f_j(X_i),
\]
the canonical morphism attached to the bifunctor
\[
 (i,j)\longmapsto f_j(X_i):I\times J\longrightarrow F.
\]
Thus Proposition~9.6 follows from the following more general assertion.
\end{prooftext}

\begin{statement}{Corollary 9.8}
Let $F$ be a $\Pi$-accessible category, let $I$ be an ordered set large with respect to $\Pi$, let $J$ be a small category such that $\card\Fl J\leq \Pi$ and such that projective limits of type $J$ are representable in $F$, and let $h:I\times J\to F$ be a functor. Then the canonical morphism
\begin{equation}
 \varinjlim_i \varprojlim_j h(i,j)
 \longrightarrow
 \varprojlim_j \varinjlim_i h(i,j)
\tag{9.8.1}
\end{equation}
is an isomorphism. In other words, the functor
\begin{equation}
 \varinjlim_i:\Hhom(I,F)\longrightarrow F
\tag{9.8.2}
\end{equation}
commutes with projective limits of type $J$, for every small category $J$ such that $\card\Fl J\leq \Pi$ and such that projective limits of type $J$ are representable in $F$. Equivalently, for every $J$ as above, the functor
\begin{equation}
 \varinjlim_j:\Hhom(J,F)\longrightarrow F
\tag{9.8.3}
\end{equation}
is $\Pi$-accessible.
\end{statement}

Since, by hypothesis, $F$ admits a conservative family of covariant representable $\Pi$-accessible functors $g:F\to(\U\text{-}\Set)$, one is immediately reduced to proving 9.8 in the case $F=\U\text{-}\Set$. We shall then prove the assertion in the form that (9.8.2) commutes with projective limits of type $J$, with $\card\Fl J\leq \Pi$. Since $I$ is filtered, (9.8.2) is known to commute with finite projective limits (2.8). A standard argument (cf. 2.3) therefore reduces us to proving that it commutes with products indexed by a set $J$ of cardinal $\leq\Pi$. This reduces to proving the bijectivity of (9.8.1) when $J$ is discrete.

For injectivity, consider two elements $a,b$ of the left-hand side. They come from $\prod_j h(i_0,j)$ for a suitable $i_0\in I$, say from elements $\prod_j a(i_0,j)$ and $\prod_j b(i_0,j)$. Suppose that the elements $\prod_j h(\infty,j)$ of the right-hand side which they define are equal; that is, for every $j$ there exists $i(j)\geq i_0$ such that $a(i_0,j)$ and $b(i_0,j)$ have the same image in $h(i,j)$. Since $I$ is large with respect to $\card J$, there exists a common upper bound $i_1\in I$ for all the $i(j)$. It follows that the two elements under consideration in $\prod_j h(i_0,j)$ have the same image in $\prod_j h(i_1,j)$, and therefore define the same element of the left-hand side of (9.8.1). This proves injectivity.

For surjectivity, let $\prod_j a(\infty,j)$ be an element of the right-hand side. For each $j\in J$, $a(\infty,j)$ comes from an element of $h(i(j),j)$ for a suitable $i(j)\in I$. As before, one can find a common upper bound $i\in I$ for all the $i(j)$. Then $\prod_j a(\infty,j)$ comes from an element $\prod_j a(i,j)$ of $\prod_j h(i,j)$, hence lies in the image of (9.8.1). This completes the proof.

\begin{statement}{Corollary 9.9}
Let $F$ be a $\Pi$-accessible (respectively accessible) category. Then, for every category $J$ such that $\card\Fl J\leq \Pi$ (respectively every small category $J$), and for every functor $(X_j)_{j\in J}:J\to F$ whose inductive limit in $F$ is representable, if the $X_j$ are $\Pi$-accessible (respectively accessible), then so is $\varprojlim X_j$.
\end{statement}

Indeed, the functor $F\to(\U\text{-}\Set)$ represented by $\varinjlim X_j$ is the projective limit of the functors represented by the $X_j$, and one applies Proposition~9.6(ii) to the resulting projective system of functors.

\begin{remarktext}{Remark 9.10}
In the statements 9.6, 9.7, 9.8, and 9.9 one may everywhere replace the words ``$\Pi$-accessible'', ``accessible'' by ``$\Pi$-preaccessible'', ``preaccessible''. The proof just given proves this variant of the preceding statements as well.
\end{remarktext}

\begin{statement}{Proposition 9.11}
Let $E$ be a category satisfying condition $L$ (9.1), and let $C$ be a small full generating subcategory (7.1). Suppose that the following condition is satisfied:
\begin{enumerate}[label=\alph*)]
\item $E$ is stable under kernels of double arrows, and the functor $\Ker$ on the category of double arrows of $E$ is accessible; for example, it commutes with small filtered inductive limits.
\end{enumerate}
Then every object of $E$ is preaccessible (9.3), and \emph{a fortiori} $E$ is preaccessible. Suppose moreover that $C$ is generating by strict epimorphisms (7.1), and suppose that $E$ satisfies the following conditions:
\begin{enumerate}[label=\alph*),resume]
\item $E$ is stable under fibre products.
\item Every small strict epimorphic family in $E$ is universally strict epimorphic (10.3); or else small coproducts are representable in $E$ and every strict epimorphism in $E$ is a universal strict epimorphism.
\end{enumerate}
Then every object of $E$ is accessible, and \emph{a fortiori} $E$ is accessible.
\end{statement}

The fact that, under (a), every object of $E$ is preaccessible follows from the fact that, for every object $X$ of $E$, the set of strict subobjects of $X$ is small (7.4), and from the following lemma.

\begin{statement}{Lemma 9.11.1}
Under the hypotheses of 9.11(a), let $X\in\ob E$ and let $\Pi$ be a cardinal such that $E$ satisfies $L_\Pi$ (9.1), the functor $\Ker$ in 9.11(a) is $\Pi$-accessible, and the set of strict subobjects of $X$ has cardinal bounded by $\Pi$. Then $X$ is $\Pi$-preaccessible.
\end{statement}

Indeed, let $I$ be a set large with respect to $\Pi$, let $(Y_i)_{i\in I}$ be an inductive system in $E$ with inductive limit $Y$, and let
\[
 (u_i,v_i:X\rightrightarrows Y_i)
\]
be a double arrow such that the composite double arrow $u,v:X\rightrightarrows Y$ satisfies $u=v$. We prove that there exists $j\geq i$ in $I$ such that $u_j=v_j$. For this, consider the inductive system of double arrows $(u_j,v_j:X\rightrightarrows Y_j)_{j\geq i}$, whose inductive limit is $(u,v)$. By hypothesis,
\[
 X=\Ker(u,v)=\varinjlim_j \Ker(u_j,v_j)=\varinjlim_j X_j,
 \qquad X_j=\Ker(u_j,v_j).
\]
The $X_j$ are strict subobjects of $X$, so there exists a subset $I'$ of $I$, consisting of indices $j\geq i$, with $\card I'\leq\Pi$, such that every $X_j$ is equal to one of the $X_{i'}$ for $i'\in I'$. Since $I$ is large with respect to $\Pi$, there is an upper bound $j$ of $I'$ in $I$. Then $X_j$ contains all the $X_{j'}$ for $j'\geq i$, hence $X_j\to X$ is an epimorphism, since the family of maps $X_{j'}\to X$ is epimorphic. It is therefore an isomorphism, since it is a strict monomorphism. Thus $u_j=v_j$, proving 9.11.1.

The second assertion of 9.11 follows from the first and from the following lemma.

\begin{statement}{Lemma 9.11.2}
Under the hypotheses of 9.11(a), (b), and (c), let $X$ be an object of $E$, and let $\Pi$ be an infinite cardinal such that $E$ satisfies $L_\Pi$ (9.1), such that $\card\ob C_{/X}\leq\Pi$, such that for two objects $X'\to X$ and $X''\to X$ of $C_{/X}$ the fibre product $X'\times_X X''$ is $\Pi$-preaccessible, and finally such that $X$ is $\Pi$-preaccessible. Then $X$ is $\Pi$-accessible.
\end{statement}

With the notation of the proof of 9.11.1, it remains only to prove that every morphism $u:X\to Y$ comes from a morphism $u_i:X\to Y_i$. Consider the family of morphisms $Y_i\to Y$, which is a strict epimorphic family. By (b) and (c), the family
\[
 Y_i\times_Y X\longrightarrow X
\]
is also strictly epimorphic. On the other hand, for every $i$, since $C$ is generating by strict epimorphisms, the family of maps
\[
 T\longrightarrow Y_i\times_Y X
\]
with source in $C$ is strictly epimorphic. In the second alternative in (c), one can find such a strict epimorphism with source a small coproduct of objects of $C$. By transitivity of universally strict epimorphic families (II~2.5), it follows that the family of maps $T_\alpha\xrightarrow{f_\alpha}X$ with source in $C$ which factor through one of the $Y_i\times_Y X$ is strictly epimorphic. Let $J$ be the set of indices of this family; its cardinal is bounded by $\card\ob C_{/X}\leq\Pi$.

For each $\alpha\in J$, choose an $i=i(\alpha)\in I$ and an $X$-morphism
\[
 T_\alpha\longrightarrow Y_i\times_Y X,
\]
or equivalently a morphism $v_\alpha:T_\alpha\to Y_i$ such that $p_i v_\alpha=uf_\alpha$, where $p_i:Y_i\to Y$ is the canonical morphism. Since $I$ is large with respect to $\Pi$, one may choose $i(\alpha)$ independently of $\alpha$; call it $i$. For every pair of indices $\alpha,\beta\in J$, consider the two composites
\[
 \operatorname{pr}_1v_\alpha,
 \operatorname{pr}_2v_\beta:
 T_\alpha\times_XT_\beta\rightrightarrows Y_i .
\]
Their composites with $Y_i\to Y$ are equal. Since $T_\alpha\times_XT_\beta$ is $\Pi$-preaccessible by hypothesis, there exists an index $i'=i(\alpha,\beta)\geq i$ such that the composites of the arrows under consideration with $Y_i\to Y_{i'}$ are equal. The set of pairs $(\alpha,\beta)$ has cardinal $\leq\Pi^2=\Pi$, because $\Pi$ is infinite; hence one may again choose $i'$ independently of $\alpha$ and $\beta$. Plainly we may suppose $i'=i$. But then, since the family $(f_\alpha:T_\alpha\to X)$ is strictly epimorphic, one can find a morphism $u_i:X\to Y_i$ such that $u_if_\alpha=v_\alpha$. We then have $p_iu_i=u$, because for every $\alpha$,
\[
 (p_iu_i)f_\alpha=p_i(u_if_\alpha)=p_iv_\alpha=uf_\alpha,
\]
and the family of the $f_\alpha$ is epimorphic. This completes the proof of 9.11.2.

\begin{remarktext}{Remark 9.11.3}
As a special case of 9.11, one sees that every object of $E$ is accessible in each of the following two cases. First, $E$ is an abelian $\U$-category with exact filtered inductive limits and admitting a small generating subcategory; here it is the second alternative in (c) which applies. Second, $E$ is the category of sheaves of sets on a topological space $X\in\U$. More generally, it is enough that $E$ be the category of sheaves of sets on a $\U$-site (II~2.1), or that $E$ be a $\U$-topos (IV~1.1).
\end{remarktext}

\begin{statement}{Definition 9.12}
Let $E$ be a category. A \emph{cardinal filtration} of $E$ is an increasing filtration $(\Filt^\Pi(E))_{\Pi\geq\Pi_0}$ of $E$ by strictly full subcategories $\Filt^\Pi(E)$, indexed by the cardinals $\Pi\in\U$ with $\Pi\geq\Pi_0$, where $\Pi_0$ is a fixed infinite cardinal depending on the cardinal filtration under consideration, and satisfying the following conditions:
\begin{enumerate}[label=\alph*)]
\item For every $\Pi\geq\Pi_0$, $\Filt^\Pi(E)$ is equivalent to a small category.

\item $E$ satisfies $L_{\Pi_0}$ (9.3), and for every $\Pi\geq\Pi_0$, $\Filt^\Pi(E)$ is stable in $E$ under filtered inductive limits indexed by ordered sets $I$ which are large with respect to $\Pi_0$ and such that $\card I\leq\Pi$.

\item For every $\Pi\geq\Pi_0$ and every $X\in\ob E$, one can find an isomorphism
\[
 X\simeq \varinjlim_I X_i,
\]
where $(X_i)_{i\in I}$ is an inductive system in $E$ indexed by an ordered set $I$ large with respect to $\Pi$ (9.1), and where the $X_i$ belong to $\Filt^\Pi(E)$. Moreover, if $X\in\ob\Filt^\Pi(E)$ and $\Pi'\geq\Pi$, one can take $I$ with $\card I\leq(\Pi')^\Pi$.
\end{enumerate}
\end{statement}

\subsection*{9.12.1}
It follows from (b) and (c) that every $X\in\ob E$ belongs to some $\Filt^\Pi(E)$ for $\Pi$ sufficiently large.

\begin{remarktext}{Example 9.12.2}
Take $E=(\U\text{-}\Set)$, $\Pi_0=\aleph_0$, and let $\Filt^\Pi(E)$ be the full subcategory of $E$ consisting of the sets $X$ such that $\card X\leq\Pi$. More generally:
\end{remarktext}

\begin{statement}{Proposition 9.13}
Let $E$ be a $\U$-category. Suppose that $E$ is stable under small filtered inductive limits, under sums of two objects, and under cokernels of double arrows; that $E$ is stable under fibre products; and that epimorphic morphisms in $E$ are universal strict epimorphisms. Let $C$ be a small full subcategory of $E$ which is generating by strict epimorphisms. Let $\Pi_0$ be an infinite cardinal $\geq\card\Fl C$. For every cardinal $\Pi\geq\Pi_0$, let $\Filt^\Pi(E)$ be the strictly full subcategory of $E$ consisting of the objects $X$ of $E$ for which there exists a strict epimorphic family $(X_i\to X)_{i\in I}$ with target $X$, such that $\card I\leq\Pi$ and $X_i\in\ob C$ for every $i\in I$. Then $(\Filt^\Pi(E))_{\Pi\geq\Pi_0}$ is a cardinal filtration of $E$. Moreover, for every $X\in\ob\Filt^\Pi(E)$, the cardinal of the set of arrows of $C_{/X}$, and \emph{a fortiori} of the set of objects of $C_{/X}$, is bounded by $\Pi^{\Pi_0}$.
\end{statement}

The last assertion is precisely 7.6.

To prove that this is a cardinal filtration, it remains to prove conditions (a), (b), and (c) of 9.12. Condition (a) follows at once from 7.5.2. For (b), suppose that
\[
 X=\varinjlim_I X_i,
\]
with $\card I\leq\Pi$ and $X_i\in\ob\Filt^\Pi(E)$. The family of canonical morphisms $X_i\to X$ is therefore strictly epimorphic. By the hypothesis on the $X_i$, for each $i\in I$ there is a strict epimorphic family $(X_{ij}\to X_i)_{j\in J_i}$ with $\card J_i\leq\Pi$. Passing to the sums $\coprod_i X_i$ and $\coprod_{i,j}X_{ij}$, one sees, by transitivity of universal strict epimorphisms (II~2.5), that the family of composites
\[
 X_{ij}\longrightarrow X_i\longrightarrow X
\]
indexed by the disjoint union of the $J_i$, for $i\in I$, is strictly epimorphic. Since $\card J\leq\Pi^2=\Pi$, the conclusion follows. Notice that we have used only the existence of $\varinjlim_I X_i$, and not the fact that $I$ is large with respect to $\Pi_0$, nor even that it is filtered.

It remains to prove (c). First note that for every $X\in\ob E$, the family of arrows $X_i\to X$ with source in $\ob C$ is strictly epimorphic, because $C$ is generating by strict epimorphisms, and is plainly small. Thus there is a cardinal $\Pi\geq\Pi_0$ such that $X\in\ob\Filt^\Pi(E)$. It remains to prove that, if $\Pi'\geq\Pi\geq\Pi_0$ and $X\in\ob\Filt^{\Pi'}(E)$, then there is an isomorphism
\[
 X\simeq \varinjlim_I X_i,
\]
where $I$ is large with respect to $\Pi$, $\card I\leq(\Pi')^\Pi$, and the $X_i$ lie in $\ob\Filt^\Pi(E)$.

Since $C$ is generating by strict epimorphisms, one has
\begin{equation*}
 X\simeq \varinjlim_{C_{/X}} T .\tag{$\ast$}
\end{equation*}
Furthermore, by successively adjoining to $C$ sums and cokernels of double arrows in $C$, and then doing the same for the category so obtained, and so on, we are reduced to the case where $C$ is stable under sums of two objects in $E$ and under cokernels of double arrows in $E$, without destroying the hypothesis $\Pi_0\geq\card\Fl C$. We may also suppose that, if $E$ contains a fixed initial object $\emptyE$, then $\emptyE\in\ob C$. This implies that the category $C_{/X}$ is stable under sums of two objects and under cokernels of double arrows. It is then filtered: it is non-empty, since otherwise the relation $(\ast)$ would show that $X$ is an initial object, so that $C_{/X}$ would contain the arrow $\emptyE\to X$, a contradiction.

Let $I$ be the set, ordered by inclusion, of full subcategories $i$ of $C_{/X}$ which are filtered and such that $\card\ob i\leq\Pi$. For every $i\in I$, let $X_i\in\ob E$ be the inductive limit of the composite functor
\[
 i\longrightarrow C_{/X}\longrightarrow E,
\]
which exists by hypothesis. Then one plainly has
\[
 \varinjlim_I X_i\simeq \varinjlim_{C_{/X}}T\simeq X.
\]
On the other hand, we have already noted that
\[
 \card\ob C_{/X}\leq(\Pi')^{\Pi_0}.
\]
It follows that $I$ is large with respect to $\Pi$ and that
\[
 \card I\leq ((\Pi')^{\Pi_0})^\Pi=(\Pi')^{\Pi\Pi_0}=(\Pi')^\Pi
\]
by the following lemma, whose proof is left to the reader; in applying it one takes $J=C_{/X}$ and $c=(\Pi')^{\Pi_0}$.

\begin{statement}{Lemma 9.13.1}
Let $J$ be a filtered category, and let $c$ and $\Pi$ be two cardinals such that
\[
 c\geq\Sup(\card\Fl J,\Pi)
\]
and such that $\card\Hom(j,j')\leq\Pi_0$ for every pair of objects $j,j'$ of $J$. Let $I$ be the set of full filtered subcategories $i$ of $J$ such that $\card\ob i\leq\Pi$. Then, ordered by inclusion, $I$ is large with respect to $\Pi$, and one has $\card I\leq c^\Pi$.
\end{statement}

This completes the proof of Proposition~9.13.

\end{document}
